\subsection{Recovering the actual twist stabilizer and faithful restrictions}

\begin{result}{E.47}{Proved}{Nodal reconstruction and faithful individual restrictions}
For every stable surface $S=S_{g,n}$ and every nonempty simplex
$\sigma$ of its procongruence curve complex, let
$\Delta_\sigma\subset U=\check\Gamma_{g,[n]}$ be its full
closed multitwist group and let $C_\sigma$ be its actual
simplex stabilizer.  Then
\[
 N_U(\Delta_\sigma)=C_\sigma.
                                                               \tag{E.240}
\]
This equality uses the identification of curves with their
closed cyclic surface subgroups and nodal reconstruction;
it does not use the disputed assertion D.26.

For every $g\geq3$ and $n\geq0$, both maps (E.230) and
(E.231) are continuous injections with closed image.
The same is true for the corresponding maps to the outer
automorphism groups of the full profinite completions.
Thus the remaining individual-group question in E.45 is
surjectivity of the displayed procongruence maps.
\end{result}
\begin{proof}
\textbf{The exact comparison for the normalizer.}
First take $\sigma$ to be a discrete multicurve, and pinch
its curves to obtain a stable nodal surface.  The vertex
groups are the full profinite fundamental groups of the
punctured cut components, and every node has its two
specified peripheral embeddings.  Denote this semi-graph
of profinite groups by $\mathcal G_\sigma$.  Van Kampen
followed by the logarithmic specialization comparison
identifies its profinite fundamental group with
$\Pi=\widehat\pi_1(S)$; each nodal cyclic group is precisely
the closure of the subgroup generated by its curve.
These identifications are up to the usual inner choices,
which are the same choices used in $\operatorname{Out}(\Pi)$.

The comparison of actual annular twists with the nodal
Dehn coordinates in
\cite[Theorem 4.8(iv), Theorem 5.7]{ET_HMTopicsI2012}
is the homomorphism
\[
 \prod_{e\in\sigma}\widehat{\mathbb Z}
       \longrightarrow\operatorname{Dehn}(\mathcal G_\sigma),
 \qquad (a_e)_e\longmapsto\prod_{e\in\sigma}T_e^{a_e}.
                                                               \tag{E.241}
\]
It is an isomorphism: Theorem 4.8 identifies the target
with the product of the nodal cyclotomes, and the local
model $xy=t_e$, with its positive meridian, sends the
$e$-th standard generator to that factor's generator.
Theorem 5.7 identifies these local generators with the
actual topological Dehn twists.  Consequently the image
in $\operatorname{Out}(\Pi)$ is exactly $\Delta_\sigma$.

All elements of $U$ preserve the set of peripheral
subgroups of $\Pi$.  In this cusp-preserving outer group,
\cite[Theorem 5.14(iii), pp.~119--121]{ET_HMTopicsI2012}
gives the exact equality
\[
 N_{\operatorname{Out}^{C}(\Pi)}(\Delta_\sigma)
          =\operatorname{Aut}(\mathcal G_\sigma).
                                                               \tag{E.242}
\]
Here the right-hand side means outer automorphisms which
preserve the actual vertex, edge and peripheral embeddings.
To record the source reduction, a scalar normalizer of a
positive nodal multitwist preserves this semi-graph by
\cite[Remark 4.2.1]{ET_HMNodal2011} and
\cite[Corollary 5.9(iii)]{ET_HMTopicsI2012}.
The scalar commensurator of the full nodal twist group is
therefore an open subgroup of
$\operatorname{Aut}(\mathcal G_\sigma)$.
An arbitrary commensurator of the full twist group
normalizes that scalar commensurator.  Commensurable
terminality of the graphic group, proved in
\cite[Corollary 2.7(iv)]{ET_Mochizuki2007CGC}, then puts it
in $\operatorname{Aut}(\mathcal G_\sigma)$.
This is the full argument of Theorem 5.14(ii)--(iii);
it also covers $n=0$.  In the proper case the graphicity
step uses the invariant monodromy filtration of
Corollary 2.7(iii) of the same 2007 paper.
Thus no choice of a nonzero cusp, nor a finite-subgroup
centralizer assertion, is required for (E.242).

A graphic element preserves the conjugacy classes of the
nodal cyclic subgroups.  By the two explicit isomorphisms
$\check C(S)\to L(\Pi)\to L'(\Pi)$ in
\cite[Theorem 4.2]{ET_Boggi2014}, intersecting the right
side of (E.242) with $U$ gives a subgroup of the actual
simplex stabilizer $C_\sigma$.
Conversely every element of $C_\sigma$ is a limit of
actual discrete stabilizers, by
\cite[Theorem 4.5]{ET_Boggi2014}.  Their annular formulas
permute the twists of $\sigma$, and hence every such
limit normalizes $\Delta_\sigma$.  This proves (E.240).
Every profinite simplex is a $U$-translate of a discrete
one: there are finitely many discrete topological
orbits, whose compact completed orbits cover the completed
complex.  Conjugating all the displayed maps proves
(E.240) for those simplices as well.
Theorem 4.2 and the proof of Theorem 4.5 precede and do
not invoke the Section 5 assertion refuted in D.26.

\textbf{The rational cut quotient detects the parameter.}
Now let $g\geq3$, and choose $g$ disjoint nonseparating
curves whose simultaneous complement is connected.
Its type is $(0,m)$ with $m=n+2g\geq6$.
The branch-pinned subgroup $C_\sigma^{\rm pin}$ of the
actual stabilizer has the exact cut sequence
\[
 1\longrightarrow\Delta_\sigma
  \longrightarrow C_\sigma^{\rm pin}
  \longrightarrow\check\Gamma_{0,m}
  \longrightarrow1.
                                                               \tag{E.243}
\]
The quotient is the full profinite pure sphere group,
by D.19.  The quotient of the whole stabilizer by
$\Delta_\sigma$ is its actual finite extension by the
permutations of cut branches and original punctures
which extend over the paired seams.  In particular,
conjugation by an element of $C_\sigma$ on the pure
quotient of (E.243) has, in its outer automorphism group,
the class of that actual permutation in
$\mathfrak S_m\subset
\operatorname{Out}(\widehat\Gamma_{0,m})$.
This follows first from cutting a discrete stabilizer;
continuity gives the same statement for its closure.
No splitting of that finite extension is used.

Take $A\in\widehat{\operatorname{GT}}$ in the kernel of
$r_{g,n}$, and choose its tower representative $F$.
Its open-object action is $\alpha=\operatorname{Inn}(h)$
for some $h\in U$.
Use the full-smoothing crossing $j_\sigma=[1]$ from
$\sigma$ to the open object.  Write
$F(j_\sigma)=u j_\sigma$ using the actual ambient
transporter coordinates of D.33.
The stabilizer of $j_\sigma$ is exactly
$\Delta_\sigma$, and hence
\[
 h\Delta_\sigma h^{-1}=u\Delta_\sigma u^{-1},
 \qquad k=u^{-1}h\in N_U(\Delta_\sigma)=C_\sigma.
                                                               \tag{E.244}
\]
For $c\in C_\sigma$, write $\bar c$ for its boundary
automorphism modulo $\Delta_\sigma$.
The actual crossing identity $j_\sigma\bar c=cj_\sigma$
gives after applying $F$
\[
             F(\bar c)=\overline{kck^{-1}}.
                                                               \tag{E.245}
\]
For completeness, this equality is read in the faithful
ambient coset $U/\Delta_\sigma$: cancel the invertible
left factor $u$ in
$u j_\sigma F(\bar c)=hch^{-1}u j_\sigma$.
Thus it is an equality of the specified boundary
quotient, rather than an assumed injectivity of a
completion of an unrelated Hom-set.

Compatibility with the pinned normalization identifies
the restriction of (E.245) to (E.243) with the standard
genus-zero action of $A$, up to its permitted inner
basepoint change.  The other side is the finite
permutation induced by $k$.
By the individually checked sphere-group comparison
\cite[Corollary 2.8(i)]{ET_HMM2022},
\[
 \operatorname{Out}(\widehat\Gamma_{0,m})
            =\mathfrak S_m\times\widehat{\operatorname{GT}},
 \qquad
       [(E.245)]=(s(k),1)=(1,A).
                                                               \tag{E.246}
\]
The displayed comparison specifies both factors: the
first is the relabeling action, and the second is the
standard labelled genus-zero action already identified
in E.16 and E.41.  Equality in this actual direct
product gives $A=1$.  This proves injectivity of
(E.230).

\textbf{Pure groups and full completions.}
Put $P=\check\Gamma_{g,n}$ and $U=\check\Gamma_{g,[n]}$.
The geometric-monodromy centralizer calculation
\cite[Theorem 6.13(ii)]{ET_HMTopicsI2012}, applied with
all primes and genus $g\geq3$, gives
\[
                   C_U(P)=1.
                                                               \tag{E.247}
\]
Indeed both groups are their actual monodromy images
inside $\operatorname{Out}^{C}(\Pi)$, and that theorem's
centralizer is the automorphism group of the universal
pointed curve, which is trivial in these types.

Suppose $(s,A)$ maps to the identity in (E.231).
Choose a lift $v\in U$ of $s$ and the compatible
automorphisms $\alpha_U,\alpha_P$ induced by $A$.
There is $p\in P$ such that
$\operatorname{Inn}(v)\alpha_P=\operatorname{Inn}(p)$.
Thus $\beta=\operatorname{Inn}(p^{-1}v)\alpha_U$
restricts to the identity on $P$.
For $u\in U$, comparison of conjugation on $P$ gives
$\beta(u)u^{-1}\in C_U(P)=1$, so $\beta=1$.
Consequently $r_{g,n}(A)=1$, whence $A=1$ by the
preceding paragraph.  Now $\ell(s)=1$; the same
centralizer equality applied to its lift implies
$s=1$.  This proves injectivity of (E.231).

The full-completion maps descend to these maps by D.31,
with the same GT parameter and the same relabeling
permutation.  A class which is inner before quotienting
is inner afterwards.  Their kernels are therefore also
trivial.  Finally all the source groups are profinite
and all targets are Hausdorff profinite outer
automorphism groups of topologically finitely generated
profinite groups.  The continuous injections are
homeomorphisms onto closed images.
\end{proof}
